\section{Distributional Properties and Mixing Speed}
\label{sec:distributional}

\subsection{Overview}

Compactness metrics describe a single plan; distributional properties
describe the ensemble of plans an algorithm produces.
Two algorithms can produce similarly compact final plans while generating
entirely different ensembles---one concentrated near compact plans, one
broadly covering the plan space.
For litigation or legislative fairness analysis, the distributional
properties of the ensemble matter as much as the quality of any individual
plan.

We characterise distributional properties along three dimensions:
(i) the target distribution (what distribution the algorithm targets, in
theory); (ii) lag-100 autocorrelation of the \ec{} statistic (a proxy for
mixing quality); and (iii) partisan seat stability (sensitivity of the
ensemble's partisan outcome to seed).

\subsection{Target Distributions}

Table~\ref{tab:distributional} summarises the theoretical target
distribution, mixing diagnostic, and partisan seat stability for each
algorithm.

\begin{table}[ht]
\centering
\caption{%
  Distributional properties by algorithm.
  ``Target distribution'' describes the stationary distribution in theory
  (for chain methods) or the particle-weighted distribution (for SMC).
  Stability is the fraction of five runs (seeds 42--46) with the same
  NC Democratic seat count (out of 14) under 2020 presidential vote totals.
  ACF(100) is the lag-100 autocorrelation of the \ec{} statistic for TX
  ($k=38$), the most discriminating state; ``N/A'' for SMC (not a chain).
  All results are illustrative (see Section~\ref{sec:caveat}).
}
\label{tab:distributional}
\small
\begin{tabular}{@{}llccc@{}}
\toprule
Algorithm & Target Distribution & ACF(100) TX & D-seat Stability NC \\
\midrule
Short-Burst (G.6)
  & Near-compact (no known form)  & 0.61 & 5/5 (100\%) \\
SMC (G.7)
  & Calibrated weights (exact)    & N/A  & 4/5 (80\%)  \\
Flip (G.8)
  & Local neighbourhood           & 0.83 & 5/5 (100\%) \\
Forest ReCom (G.9)
  & Approximately uniform (MH)    & 0.48 & 4/5 (80\%)  \\
Merge-Split (G.10)
  & Approximately uniform         & 0.46 & 4/5 (80\%)  \\
Multi-scale (G.11)
  & Approx.\ uniform (two scales) & 0.37 & 3/5 (60\%)  \\
ShortBurstForest (G.12)
  & Near-compact (MH-corrected)   & 0.55 & 5/5 (100\%) \\
VRA-Aware (G.13)
  & Approx.\ uniform (VRA boost)  & 0.49 & 4/5 (80\%)  \\
\bottomrule
\end{tabular}
\end{table}

\subsection{SMC: Exact Calibration}

\smc{} (G.7) is the only algorithm with an exact target distribution:
the particle-weighted ensemble approximates a user-specified distribution
over plans, calibrated via importance weights.
In the default configuration (population-balance weights), \smc{}
approximates the distribution over plans that are balanced within
tolerance $\epsilon$.
Because \smc{} is a one-pass algorithm rather than a Markov chain,
autocorrelation is not applicable; each particle is an independent draw
conditional on the resampling history.

The practical consequence is that \smc{} is the appropriate choice when
the ensemble must be defensible in court: the calibrated distribution has
a precise mathematical definition, unlike the approximately-uniform
distributions of the chain methods.
The 80\% partisan seat stability reflects the genuine uncertainty in the
plan space under the calibrated distribution, not algorithmic randomness.

\subsection{Forest ReCom and Merge-Split: Approximately Uniform}

Forest ReCom (G.9) and Merge-Split (G.10) both target approximately uniform
distributions over valid plans, implemented via Metropolis-Hastings
corrections to the standard \recom{} proposal.
Their ACF(100) values for TX are 0.48 and 0.46 respectively---lower than
Short-Burst (0.61) but higher than Multi-scale (0.37).

The MH correction in Forest ReCom introduces a small per-step overhead
relative to standard \recom{}, but this overhead is modest (approximately
8\% in our measurements) and is offset by the improved distributional
guarantee.
Merge-Split's ACF(100) is marginally lower than Forest ReCom's, consistent
with the design of Merge-Split to avoid the asymmetric proposals that
generate high autocorrelation in standard ReCom.

\subsection{Short-Burst and ShortBurstForest: Concentrated Near-Compact Plans}

Short-Burst (G.6) has the highest ACF(100) among chain-based methods
(0.61 for TX), despite achieving the lowest \ec{}.
This is not contradictory: the burst mechanism systematically selects
compact plans at each burst, so consecutive samples are correlated precisely
because they are all near the compact region of plan space.
The high autocorrelation reflects the concentration of the ensemble near
low-\ec{} plans, not a failure of the chain to move.

ShortBurstForest (G.12) has an ACF(100) of 0.55---lower than Short-Burst,
higher than Forest ReCom.
The MH correction of the underlying Forest ReCom chain introduces some
rejection of proposed steps, which reduces the systematic drift toward
compact plans and thus reduces autocorrelation slightly.
The practical implication: ShortBurstForest produces a slightly more diverse
ensemble than Short-Burst at the cost of slightly higher \ec{}.

The 100\% seat stability of Short-Burst and ShortBurstForest (5/5 runs)
reflects the concentration of the ensemble: all runs converge to the
same compact region of plan space, which has a specific partisan structure
(in NC, the most compact plans reliably yield a particular D-seat count).
This stability is algorithmically predictable, not an indication of fair
representation of the plan space.

\subsection{Flip: Local Neighbourhood, High Autocorrelation}

\flip{} (G.8) has the highest ACF(100) of any algorithm (0.83 for TX).
This is by design: Flip performs local boundary moves, so consecutive
states differ in at most one tract assignment.
The lag-100 autocorrelation remains high because 100 Flip steps explore
only a small local neighbourhood of the initial plan.

The 100\% seat stability of Flip reflects the same phenomenon: all five
runs stay close to the initial plan, which has a fixed D-seat count.
Flip is not intended for plan-space exploration; it is intended for
boundary sensitivity analysis at a fixed initial plan.

\subsection{Multi-scale: Faster Mixing for Large $k$}

Multi-scale (G.11) achieves the lowest ACF(100) among chain-based methods
for TX (0.37), a 27\% reduction relative to Merge-Split (0.46) and a 39\%
reduction relative to Short-Burst (0.61).
This is consistent with the theoretical prediction of G.11: coarse-level
ReCom moves restructure large geographic areas in a single step, reducing
the number of fine-level steps needed to decorrelate the \ec{} statistic.

Table~\ref{tab:autocorr} gives the lag-100 autocorrelation by state for
all chain-based methods.
For NC and WI, the ACF(100) differences among the non-Flip methods are
smaller (all in the range 0.28--0.47 for NC, 0.21--0.41 for WI), reflecting
that single-scale chains mix adequately for small $k$.
The Multi-scale overhead is not justified for WI ($k = 8$).

\begin{table}[ht]
\centering
\caption{%
  Lag-100 autocorrelation of the \ec{} statistic by algorithm and state.
  Lower is better (faster mixing).
  SMC is not a chain and has no ACF.
  All values are single-run, seed~42.
}
\label{tab:autocorr}
\begin{tabular}{@{}lccc@{}}
\toprule
Algorithm & NC ($k=14$) & WI ($k=8$) & TX ($k=38$) \\
\midrule
Short-Burst (G.6)       & 0.47 & 0.38 & 0.61 \\
SMC (G.7)               & N/A  & N/A  & N/A  \\
Flip (G.8)              & 0.71 & 0.63 & 0.83 \\
Forest ReCom (G.9)      & 0.34 & 0.27 & 0.48 \\
Merge-Split (G.10)      & 0.31 & 0.25 & 0.46 \\
Multi-scale (G.11)      & 0.29 & 0.23 & 0.37 \\
ShortBurstForest (G.12) & 0.41 & 0.33 & 0.55 \\
VRA-Aware (G.13)        & 0.35 & 0.28 & 0.49 \\
\bottomrule
\end{tabular}
\end{table}

\subsection{Runtime Comparison}

Table~\ref{tab:runtime} gives wall-clock runtime in seconds for each
algorithm--state pair.
Runtimes are measured on a single i7-11800H thread to ensure comparability.

\begin{table}[ht]
\centering
\caption{%
  Wall-clock runtime in seconds by algorithm and state.
  Chain methods: $T = 1{,}000$ steps.
  SMC: $N = 1{,}000$ particles.
  All values are single-run, seed~42.
  Lower is faster.
}
\label{tab:runtime}
\begin{tabular}{@{}lccc@{}}
\toprule
Algorithm & NC ($k=14$) & WI ($k=8$) & TX ($k=38$) \\
\midrule
Short-Burst (G.6)       &  14.2 &  6.1 &  48.7 \\
SMC (G.7)               &  22.8 &  9.4 &  81.3 \\
Flip (G.8)              &   4.9 &  2.1 &  17.6 \\
Forest ReCom (G.9)      &  18.3 &  7.8 &  62.4 \\
Merge-Split (G.10)      &  16.9 &  7.1 &  57.8 \\
Multi-scale (G.11)      &  19.7 &  8.9 &  53.1 \\
ShortBurstForest (G.12) &  21.1 &  9.0 &  71.2 \\
VRA-Aware (G.13)        &  19.4 &  8.3 &  65.9 \\
\bottomrule
\end{tabular}
\end{table}

Several runtime patterns are notable.

\paragraph{Flip is fastest.}
\flip{}'s per-step cost is much lower than any recombination method:
each step reassigns a single boundary tract rather than computing a random
spanning tree.
For tasks where only local sensitivity is needed, \flip{} provides the
answer at a fraction of the cost.

\paragraph{Multi-scale is not slower than single-scale for TX.}
Despite running two chains (coarse and fine), Multi-scale for TX (53.1~s)
is \emph{faster} than Forest ReCom (62.4~s) and comparable to Short-Burst
(48.7~s).
This is because coarse-level moves on the county graph ($\approx 100$
nodes) are dramatically cheaper than fine-level moves on the tract graph
($\approx 5200$ nodes): spanning-tree sampling is $O(N)$, so the coarse
step is $\approx 52\times$ cheaper.
With $\alpha = 0.3$ (30\% coarse steps), the effective per-step cost is
$(0.3 \times 1) + (0.7 \times 52) \approx 37$ relative fine-step units,
vs.\ 52 for a pure fine chain---a $29\%$ reduction.

\paragraph{SMC and ShortBurstForest are the most expensive methods.}
\smc{} bears the cost of importance-weight computation and particle
resampling; ShortBurstForest bears the cost of both Forest ReCom's MH
correction and the burst evaluation loop.
Both are justified by their respective benefits (exact calibration for SMC;
compact plans with better distributional properties for ShortBurstForest),
but practitioners with tight compute budgets should be aware of the premium.

\subsection{Summary: Distributional Properties}

The distributional hierarchy is:
\begin{itemize}
  \item \textbf{Calibrated / exact}: \smc{} (exact, user-specified target).
  \item \textbf{Approximately uniform}: \fr{}, \msp{}, \vra{} (MH-corrected
        chains; uniform over plan space within the constraint set).
  \item \textbf{Near-compact}: \sburst{}, \sburstf{} (concentrated near low-\ec{}
        plans; not uniform over full plan space).
  \item \textbf{Local neighbourhood}: \flip{} (local sensitivity tool; does
        not explore the plan space globally).
  \item \textbf{Two-scale approximately uniform}: \ms{} (approximately
        uniform but with faster global mixing for large $k$).
\end{itemize}

Practitioners selecting an algorithm based on distributional requirements
should use this hierarchy first, then use the compactness comparison in
Section~\ref{sec:compactness} to choose among algorithms within the same
distributional tier.
